%\setcounter{secnumdepth}{+2}
\chapter{Introdução}
%\chapter{Introduction}
\label{cap1}

\section{Introdução}

O presente trabalho surge no âmbito do Mestrado em Design e Desenvolvimento de 
Jogos Digitais da \ac{UBI} e foca-se na conceção e implementação de um videojogo de 
simulação e gestão de infraestruturas tecnológicas. O projeto propõe uma 
experiência imersiva onde o jogador assume o papel de gestor de um pequeno 
centro de dados (data center), tendo como objetivo primordial a viabilidade 
económica da operação por uma gestão estratégica de recursos. A abordagem 
adotada procura equilibrar o rigor técnico de conceitos reais de computação com 
o entretenimento, utilizando métricas e componentes de hardware autênticos para 
reforçar a imersão do utilizador.

A estrutura deste relatório de projeto reflete o percurso de desenvolvimento do jogo, 
começando por contextualizar a relevância do tema e as motivações que levaram à 
sua escolha. Ao longo dos capítulos subsequentes, é apresentado o estado da arte, 
onde se analisam títulos de referência no mercado e literatura científica relevante, 
seguido do detalhe das decisões de design e das técnicas de implementação utilizadas. 
Finalmente, o documento descreve os processos de teste e avaliação que validam o protótipo 
desenvolvido, culminando numa reflexão sobre as contribuições do trabalho e as perspetivas 
para desenvolvimentos futuros na área dos jogos educativos e de simulação.

\section{Motivação}

A motivação para o desenvolvimento deste projeto surge da intersecção entre 
a complexidade da gestão de infraestruturas tecnológicas e o potencial lúdico 
dos simuladores modernos. Embora o tema de gestão de data centers possa parecer,
à primeira vista, estritamente técnico ou árido, existe um ecossistema rico 
em mecânicas de estratégia, logística e resolução de problemas que se 
traduzem naturalmente em experiências de jogo envolventes.

Um dos principais motores deste trabalho é a vontade de desmistificar 
conceitos de computação de alto desempenho, como o consumo energético em Watts 
ou a capacidade de processamento em TFLOPS, integrando-os num ciclo de jogo 
onde o conhecimento é uma ferramenta para o sucesso e não uma barreira à 
entrada. A utilização de unidades e métricas reais visa reforçar a imersão e 
conferir um peso estratégico às decisões do jogador, permitindo uma aprendizagem 
orgânica através da simulação.

Contudo, a motivação central não reside apenas na simulação fiel, mas sim na 
criação de um produto onde a diversão é o fator predominante. Ao contrário de 
simuladores puramente técnicos, este projeto adota uma estética inspirada na 
ficção e elementos de fantasia para garantir um ambiente visualmente apelativo 
e menos formal. A inclusão de sistemas de combate e mini-jogos baseados em 
terminais virtuais serve o propósito vital de quebrar a monotonia da gestão 
económica, introduzindo dinamismo e momentos de adrenalina que mantêm o jogador 
num estado de flow constante.

Em suma, este projeto é motivado pela convicção de que é possível construir um 
jogo "sério" que não se leve demasiado a sério. O objetivo é proporcionar uma 
experiência recompensadora onde o jogador se sinta instigado a otimizar o seu 
centro de dados, não por obrigação académica, mas pelo puro entretenimento de 
superar desafios, gerir recursos e proteger a sua infraestrutura num universo 
digital vibrante.

\subsection{Organização do Documento}

O presente documento encontra-se estruturado em seis capítulos, os quais refletem as diferentes etapas do ciclo de desenvolvimento do projeto, desde a sua concetualização e fundamentação teórica até à implementação, testagem e reflexão crítica dos resultados obtidos. A organização dos capítulos apresenta-se da seguinte forma:

\begin{enumerate}
    \item \textbf{Introdução}: Este capítulo contextualiza o enquadramento geral do trabalho, apresentando a temática da simulação e gestão de \textit{data centers}. São detalhadas as motivações subjacentes ao projeto, os objetivos a atingir, a metodologia adotada e a estrutura formal deste relatório, acompanhada pela respetiva calendarização das tarefas.

    \item \textbf{Estado da Arte}: Apresenta uma análise crítica do panorama atual, dividida entre o estudo de videojogos comerciais de referência na área da simulação, montagem e gestão tecnológica (\textit{e.g.}, \textit{Data Center Demo}, \textit{PC Building Simulator}, \textit{Startup Company}, \textit{Grey Hack} e \textit{Software Inc.}) e a revisão de literatura científica relevante sobre jogos sérios, aprendizagem indireta e mecânicas de jogo.

    \item \textbf{Técnicas e Decisões de Design}: Descreve a fundamentação teórica do \textit{game design}, abordando os modelos e \textit{frameworks} aplicados (\textit{MDA}, \textit{Flow} e Modelo de \textit{Hook}). Adicionalmente, detalha as decisões criativas e concetuais relativas à interface de utilizador (\textit{UI/UX}), \textit{level design}, sistemas de notificação e correio eletrónico, referências estéticas para os elementos 3D, bem como as mecânicas de economia, configuração de equipamento e combate.

    \item \textbf{Implementação do Jogo}: Expõe o processo prático de desenvolvimento do videojogo, discriminando as ferramentas e tecnologias utilizadas (como o motor \textit{Unity}, o \textit{software Blender}, linguagens de programação e ferramentas de apoio com Inteligência Artificial). Inclui a engenharia de software (diagramas de arquitetura de sistema e casos de uso), a criação de modelos 3D e \textit{assets} 2D, e o desenvolvimento técnico de subsistemas fulcrais, tais como o sistema de ataques/defesa e o \textit{minigame} \textit{WordRush}.

    \item \textbf{Testes e Avaliação do Jogo}: Detalha a metodologia de teste aplicada com utilizadores reais e a consequente recolha de dados quantitativos e qualitativos. Este capítulo abrange a análise da \ac{UX} com base nas heurísticas de usabilidade de Nielsen, a avaliação do conhecimento técnico transmitido aos jogadores (através de questionários pré e pós-jogo) e o estudo dos dados analíticos/telemetria recolhidos durante as sessões de jogo.

    \item \textbf{Conclusões e Trabalho Futuro}: Sintetiza as principais conclusões retiradas ao longo do projeto, destacando o cumprimento dos objetivos propostos e as contribuições concretas do trabalho. Por fim, são perspetivadas várias linhas de desenvolvimento futuro e melhorias a implementar em edições subsequentes do jogo.
\end{enumerate}

\section{Calendarização do projeto}
\todo{Revisar}

A Tabela \ref{tab:Calendario}, apresenta a calendarização do projeto. É de notar que as tarefas apresentadas foram realizadas de forma intercalada.

\begin{table}[H]
\begin{tabular}{|l|r|}
\textbf{Tarefa} & \textbf{Duração} \\ \hline
Prototipagem rápida e busca de referências & 1 mês \\ \hline
Implementação do \textit{minigame WordRush} & 1 mês \\ \hline
Design dos sistemas & 1 mês \\ \hline
Implementação dos sistemas & 2,5 meses \\ \hline
Referências visuais de 3D e \ac{VFX} & 1 mês \\ \hline
Criação dos \ac{VFX} e dos modelos 3D & 1 mês \\ \hline
Implementação de efeitos sonos e narração & 1 semana \\ \hline
Correção de \textit{bugs} & 1 mês \\ \hline
Realização dos testes & 1 mês \\ \hline
Escrita do relatório & 2 meses \\ \hline
\end{tabular}
\centering
\caption{Calendarização do projeto}
\label{tab:Calendario}
\end{table}

\todo{Aqui}